\section{Speed Layer}
\label{sec:speed}

\subsection{Kafka Ingestion}

The speed layer uses Kafka as the real-time ingestion backbone. The topic \texttt{cybersecurity-logs} receives JSON log events produced from the historical dataset. The producer uses the source IP as the event key, which keeps related events naturally grouped for downstream analysis. Kafka offsets were checked to confirm that the stream is receiving data.

\begin{figure}[H]
    \centering
    \includegraphics[width=0.88\textwidth]{\detokenize{figures/Speed Layer- Hamza/kafkatopic.png}}
    \caption{Kafka topic verification for \texttt{cybersecurity-logs}.}
    \label{fig:kafka-topic}
\end{figure}

\subsection{Spark Streaming Detection Rules}

Spark Structured Streaming consumes the Kafka topic, parses JSON messages into a DataFrame, applies detection rules, and writes the results to Cassandra. The implemented real-time rules match the project statement:

\begin{itemize}
    \item \textbf{Brute-force detection:} 5 or more blocked/failed events from the same source IP in about one minute.
    \item \textbf{Signature detection:} suspicious tools and payloads such as \texttt{sqlmap}, \texttt{nikto}, \texttt{nmap}, SQL injection expressions, sensitive paths, and XSS-like strings.
    \item \textbf{Abnormal volume detection:} high bytes transferred by source IP in a 10-second window. The academic target is 10 MB/10 seconds; the experiment also includes a calibrated threshold from observed dataset traffic to generate demonstrable alerts.
    \item \textbf{Threat score:} aggregation by source IP using malicious count, suspicious count, brute-force indicators, signature alerts, and traffic volume.
\end{itemize}

\begin{figure}[H]
    \centering
    \includegraphics[width=0.90\textwidth]{\detokenize{figures/Speed Layer- Hamza/Screenshot 2026-05-13 194424.png}}
    \caption{Real-time signature detection and Cassandra write path.}
    \label{fig:signature-detection}
\end{figure}

\begin{figure}[H]
    \centering
    \includegraphics[width=0.82\textwidth]{\detokenize{figures/Speed Layer- Hamza/threat_score.png}}
    \caption{Real-time threat score computation per source IP.}
    \label{fig:threat-score}
\end{figure}

\subsection{Cassandra Active Threat Store}

Cassandra stores the active and recent state of the speed layer. The schema follows the access patterns of the API: recent logs by IP, active threats by IP, signature alerts, volume alerts, and aggregate scores.

\begin{table}[H]
    \centering
    \small
    \caption{Cassandra tables used by the speed layer.}
    \label{tab:cassandra-tables}
    \begin{tabularx}{\textwidth}{>{\ttfamily}p{0.27\textwidth} X}
        \toprule
        Table & Role \\
        \midrule
        logs & Raw or normalized events consumed from Kafka \\
        realtime\_threats & Active threats such as brute-force detections \\
        signature\_alerts & Alerts produced by known signatures and suspicious paths \\
        volume\_alerts & Abnormal transferred-byte windows by source IP \\
        threat\_scores & Aggregated risk score per source IP \\
        \bottomrule
    \end{tabularx}
\end{table}

\begin{figure}[H]
    \centering
    \includegraphics[width=0.72\textwidth]{\detokenize{figures/Speed Layer- Hamza/Capture Cassandra/describe_tables_cassandra.png}}
    \caption{Cassandra keyspace tables used by Spark Streaming and the API.}
    \label{fig:cassandra-tables}
\end{figure}

\subsection{Speed Layer Assessment}

The speed layer satisfies the real-time requirements because it uses Kafka for continuous ingestion, Spark Streaming for low-latency detection, and Cassandra for recent active threats. It detects at least three threat categories in real time: brute force, signature-based attacks, and abnormal volume. The additional threat score makes the output directly usable by the dashboard and the recommendation API.
